\documentclass[10pt,twocolumn]{article}

% ─── Packages ───────────────────────────────────────────────────────────────
\usepackage[margin=1in]{geometry}
\usepackage{times}
\usepackage{graphicx}
\usepackage{amsmath,amssymb,amsthm}
\usepackage{booktabs}
\usepackage{multirow}
\usepackage{hyperref}
\usepackage{xcolor}
\usepackage{algorithm}
\usepackage{algorithmic}
\usepackage{caption}
\usepackage{subcaption}
\usepackage{url}
\usepackage{cite}
\usepackage{microtype}
\usepackage{enumitem}

\hypersetup{
    colorlinks=true,
    linkcolor=blue!70!black,
    citecolor=green!50!black,
    urlcolor=blue!70!black
}

% ─── Title ──────────────────────────────────────────────────────────────────
\title{
  \textbf{Consistency-Regularised Training for Corruption Robustness:\\
  A KL-Divergence Warm-Up Approach on CIFAR-10}
}

\author{
Arshdeep Singh\\
\textit{Computer Science Engineering}\\
\textit{Thapar Institute of Engineering and Technology}\\
asingh37\_be23@thapar.edu
\\[1em]
Sukhmanpreet Singh\\
\textit{Computer Science Engineering}\\
\textit{Thapar Institute of Engineering and Technology}\\
ssingh33\_be23@thapar.edu
}

\date{}

% ─── Document ───────────────────────────────────────────────────────────────
\begin{document}
\maketitle

% ─── Abstract ───────────────────────────────────────────────────────────────
\begin{abstract}
Deep neural networks achieve high accuracy on clean test data but remain brittle when inputs are corrupted by common real-world distortions such as noise, blur, weather effects, and digital artefacts. We propose a \emph{consistency-regularised training} (CRT) scheme that directly addresses this brittleness by penalising the KL divergence between the network's softmax distributions on clean and lightly corrupted views of the same training image. A linear warm-up schedule on the KL coefficient allows the network to first acquire class-discriminative features through standard cross-entropy before the consistency pressure is gradually introduced. We evaluate CRT on CIFAR-10 and CIFAR-10-C --- the standard 15-corruption, 5-severity benchmark --- using a CIFAR-adapted ResNet-18 backbone. Our method is compared against a standard cross-entropy baseline and AugMix, a strong published robustness baseline. Experimental results show that CRT consistently improves mean corrupted accuracy over the baseline across all severity levels, with particularly pronounced gains on noise-type corruptions. We provide a comprehensive ablation study demonstrating the importance of both the KL regularisation term and the warm-up schedule. Code and checkpoints are publicly available.
\end{abstract}

% ─── 1. Introduction ────────────────────────────────────────────────────────
\section{Introduction}

Modern convolutional neural networks (CNNs) achieve near-human performance on curated image classification benchmarks~\cite{khazem2026macs,erichson2024noisymix}. However, this strong performance often degrades substantially when the same networks are evaluated under common real-world image corruptions --- blur, noise, weather effects, and digital compression artefacts~\cite{ranabhat2025shapebias,kumar2026demhec}. Such distribution shift is not adversarial in nature; it is the everyday imperfection encountered when a trained model is deployed on hardware whose sensor, lighting, or transmission conditions differ from the training environment.

The vulnerability of deep networks to these \emph{common corruptions} is well documented. Ranabhat \emph{et al.}~\cite{ranabhat2025shapebias} and Kumar \emph{et al.}~\cite{kumar2026demhec} both demonstrate that standard classifiers suffer significant accuracy degradation across a range of corruption types and severity levels, even when those architectures perform well on clean evaluation sets. This body of evidence has motivated a line of work aimed at improving \emph{corruption robustness} rather than worst-case adversarial robustness.

Data augmentation is a natural tool for improving robustness. Training on corrupted or heavily augmented images can reduce the gap between clean and corrupted performance~\cite{erichson2024noisymix,kumar2026demhec}. AugMix~\cite{hendrycks2020augmix} takes this a step further by combining multiple stochastic augmentation chains and adding a Jensen-Shannon divergence (JSD) consistency penalty between clean and augmented predictions, achieving state-of-the-art corruption robustness at the time of publication.

\paragraph{Motivation.} Despite its effectiveness, AugMix has two practical limitations. First, each training step requires three forward passes (one clean, two augmented), which increases compute per epoch. Second, the JSD penalty is applied with a fixed weight throughout training, which can destabilise early learning when the network has not yet formed reliable categorical representations.

Inspired by these observations, we propose \textbf{Consistency-Regularised Training (CRT)}, a simpler two-view method that enforces prediction consistency between a clean image and a single lightly corrupted counterpart. The objective is:

\begin{equation}
\mathcal{L} = \underbrace{\mathrm{CE}(f(x), y)}_{\text{supervised}} + \underbrace{\mathrm{CE}(f(x'), y)}_{\text{corrupted}} + \underbrace{\lambda_t \cdot D_{\mathrm{KL}}\!\left(p(x) \,\|\, p(x')\right)}_{\text{consistency}},
\label{eq:loss}
\end{equation}

where $x$ is a clean image, $x'$ is a lightly corrupted view, $p(\cdot)$ denotes the softmax output, and $\lambda_t$ is a linearly increasing KL coefficient subject to a warm-up schedule over the first $T_w$ epochs.

\paragraph{Contributions.} We make the following contributions:
\begin{enumerate}[leftmargin=*, topsep=2pt, itemsep=1pt]
  \item We propose CRT, a simple two-view consistency loss with a KL warm-up schedule for improving common-corruption robustness.
  \item We show that CRT outperforms the standard cross-entropy baseline on all 15 CIFAR-10-C corruption types.
  \item We provide a direct comparison with AugMix on CIFAR-10-C, demonstrating competitive performance with lower training complexity.
  \item We conduct ablation studies isolating the effects of the KL term and the warm-up schedule.
\end{enumerate}

% ─── 2. Related Work ────────────────────────────────────────────────────────
\section{Related Work}

\subsection{Corruption Robustness Benchmarks}

Erichson \emph{et al.}~\cite{erichson2024noisymix} and Ranabhat \emph{et al.}~\cite{ranabhat2025shapebias} both evaluate their methods against standardised corruption benchmarks such as CIFAR-10-C and ImageNet-C, which apply 15 corruption types at 5 severity levels to enable rigorous measurement of model degradation under distribution shift. A consistent finding across this literature is that architectural improvements alone --- such as deeper or wider networks --- yield negligible gains in corruption robustness relative to clean accuracy gains, underscoring the need for dedicated training-time robustness techniques. Bekdache and Shanbhag~\cite{bekdache2025federl} further extend this evaluation setting to federated learning scenarios, confirming that corruption robustness remains a challenge even when client models are trained on clean data but deployed under noisy real-world conditions.

\subsection{Data Augmentation for Robustness}

Data augmentation is among the most widely used tools for improving generalisation. Erichson \emph{et al.}~\cite{erichson2024noisymix} introduced NoisyMix, a training scheme that leverages noisy augmentations applied in both input and feature space to jointly improve model robustness and in-domain accuracy. Evaluated on ImageNet-C, ImageNet-R, and ImageNet-P, NoisyMix demonstrates that combining noise-based augmentation with implicit regularisation in the feature space yields meaningful gains in corruption robustness. While NoisyMix focuses primarily on noise-driven augmentation across multiple representation levels, it is not paired with an explicit prediction-consistency objective of the kind we adopt.

Kumar \emph{et al.}~\cite{kumar2026demhec} proposed DEM-HEC, an entropy-guided fine-tuning framework that strengthens corruption robustness while preserving clean accuracy. The key observation motivating DEM-HEC is that common corruptions tend to collapse the network's internal feature space into a high-entropy state, causing the model to rely on a small and fragile subset of features. To counter this, DEM-HEC generates high-entropy samples within a bounded perturbation region and aligns them with clean embeddings using a contrastive loss, combined with knowledge distillation from a teacher snapshot. Our method is complementary to DEM-HEC in that CRT operates on a single corrupted view generated at a controlled severity, rather than gradient-ascent-based high-entropy samples, and enforces consistency directly in the softmax output space via KL divergence rather than in the embedding space via contrastive objectives.

Hendrycks \emph{et al.}~\cite{hendrycks2020augmix} introduced AugMix, the most directly related prior work. AugMix produces two stochastic augmentation chains applied to the same image and imposes a Jensen-Shannon divergence consistency penalty between clean and augmented logit distributions, in addition to the supervised cross-entropy. This three-view approach substantially reduces Mean Corruption Error (mCE) on ImageNet-C and CIFAR-10-C. Erichson \emph{et al.}~\cite{erichson2024noisymix} build on this direction by extending noisy augmentation into the feature space, demonstrating that operating at multiple levels of representation compounds the robustness benefit. Our method differs from both AugMix and NoisyMix in three key respects: (i) we use a single corrupted view rather than two augmentation chains; (ii) we adopt an asymmetric KL divergence rather than symmetric JSD; and (iii) we introduce a warm-up schedule for the consistency coefficient, which is absent in both prior approaches.

\subsection{Consistency Regularisation}

Consistency regularisation has attracted considerable attention as a principled strategy for enforcing stable predictions under input variation. Ranabhat \emph{et al.}~\cite{ranabhat2025shapebias} proposed two complementary regularisation strategies for CNNs that promote shape-biased representations as a route to improved corruption robustness. The first introduces an auxiliary loss that enforces feature consistency between original and low-frequency filtered inputs, discouraging dependence on high-frequency textures; the second incorporates supervised contrastive learning to structure the feature space around class-consistent, shape-relevant representations. Both strategies improve CIFAR-10-C corruption robustness without degrading clean accuracy, highlighting that the \emph{type} of consistency enforced matters as much as the consistency objective itself.

Khazem~\cite{khazem2026macs} introduced MaCS (Margin and Consistency Supervision), a closely related framework that augments cross-entropy with a hinge-squared margin penalty and a KL divergence consistency regulariser between clean inputs and mildly perturbed views. MaCS is architecture-agnostic and provides a unifying theoretical analysis demonstrating that increasing classification margin while reducing local sensitivity yields improved generalisation guarantees and a provable robustness radius bound. Our work shares the KL-based consistency objective with MaCS but differs in two important ways: we introduce a \emph{linear warm-up schedule} on the KL coefficient, which we show is essential for stabilising early training, and we restrict our evaluation to the CIFAR-10-C benchmark with a controlled corruption pool rather than the general perturbation setting.

Bekdache and Shanbhag~\cite{bekdache2025federl} address corruption robustness under resource constraints via FedERL, which incorporates a data-agnostic robust training (DART) procedure that operates server-side and applies consistency-based objectives without requiring access to the original training data. Their work underscores that consistency regularisation can be decoupled from the standard training pipeline and applied post-hoc, a direction complementary to our inline KL warm-up scheme. Furthermore, Cui \emph{et al.}~\cite{cui2025gkl} provide a detailed theoretical analysis of the KL divergence loss, showing it is equivalent to a decoupled formulation comprising a weighted mean-square-error term and a soft-label cross-entropy term. Their generalised GKL loss addresses the asymmetric optimisation behaviour of standard KL divergence, a finding that motivates the asymmetric direction we adopt in our CRT objective: we penalise $D_{\mathrm{KL}}(p(x) \| p(x'))$ rather than the reverse, ensuring that the corrupted view's predictive distribution is pulled towards the clean one rather than vice versa.

% ─── 3. Methodology ─────────────────────────────────────────────────────────
\section{Methodology}

\subsection{Problem Setting}

Let $\mathcal{D} = \{(x_i, y_i)\}_{i=1}^{N}$ be the clean training set, where $x_i \in \mathbb{R}^{3 \times 32 \times 32}$ and $y_i \in \{1, \ldots, C\}$ for $C=10$ classes. We train a classifier $f_\theta : \mathbb{R}^{3 \times H \times W} \to \Delta^{C-1}$ parameterised by $\theta$, where $f_\theta(x) = \mathrm{softmax}(g_\theta(x))$ and $g_\theta$ produces logits. Our goal is to improve the classifier's accuracy not only on the clean test set but on the CIFAR-10-C corruption benchmark.

\subsection{Corrupted View Generation}

For each training image $x$, we generate a corrupted view $x'$ by applying one of five controlled corruption types, sampled uniformly: (i) Gaussian noise, (ii) Gaussian blur, (iii) JPEG compression, (iv) brightness shift, and (v) contrast adjustment. All corruptions are applied at a fixed severity of $s=2$ (on a 1--5 scale matching CIFAR-10-C) to ensure the corrupted view remains recognisable and the label $y$ is preserved. This severity was chosen to be mild enough that the corrupted view is still informative for the classification task, while being strong enough to create meaningful distribution shift for the consistency loss.

\subsection{Training Objective}

The CRT objective at epoch $t$ for a minibatch $\mathcal{B}$ is:

\begin{equation}
\begin{aligned}
\mathcal{L}_t = \frac{1}{|\mathcal{B}|} \sum_{(x,y) \in \mathcal{B}} \Bigl[ 
\mathrm{CE}(g_\theta(x), y) + \mathrm{CE}(g_\theta(x'), y) \\
+ \lambda_t \cdot D_{\mathrm{KL}}(p(x) \| p(x')) 
\Bigr]
\end{aligned}
\label{eq:full_loss}
\end{equation}

where $\mathrm{CE}(g, y) = -\log \mathrm{softmax}(g)_y$ is the cross-entropy loss and $D_{\mathrm{KL}}(p \| q) = \sum_c p_c \log(p_c / q_c)$ is the KL divergence from the clean distribution to the corrupted distribution.

\paragraph{Rationale for KL vs.\ JSD.} We use the asymmetric KL divergence $D_{\mathrm{KL}}(p(x) \| p(x'))$ rather than the symmetric JSD used in AugMix. This choice enforces that the corrupted view's predictive distribution closely tracks the clean view, which is the semantically correct direction: the ground truth resides in the clean signal, and the corrupted view should not diverge from it. This asymmetry is theoretically well-motivated: Cui \emph{et al.}~\cite{cui2025gkl} demonstrate that standard KL divergence exhibits asymmetric gradient behaviour, where optimisation pressure differs meaningfully depending on the direction of the divergence. Empirically, we find this asymmetric formulation more stable in early training, consistent with their analysis of convergence challenges under symmetric or reversed KL objectives.

\subsection{KL Warm-Up Schedule}

We introduce a linear warm-up schedule for $\lambda_t$:

\begin{equation}
\lambda_t = \lambda_{\max} \cdot \min\!\left(1, \frac{t}{T_w}\right),
\label{eq:warmup}
\end{equation}

where $\lambda_{\max} = 0.5$ is the maximum KL weight and $T_w = 10$ is the warm-up duration in epochs. For $t < T_w$, the network trains primarily with cross-entropy on both clean and corrupted views, with the consistency pressure ramping up smoothly. The schedule is motivated by the observation that in early training, the softmax distributions $p(x)$ and $p(x')$ are highly uncertain and may not meaningfully reflect class structure; applying a strong consistency loss at this stage can impede feature learning. Figure~\ref{fig:warmup} illustrates the warm-up behavior, showing how the KL loss and weight evolve during training.

\begin{figure}[t]
\centering
\includegraphics[width=0.95\columnwidth]{consistency_warmup.png}
\caption{KL warm-up behavior during training. The KL weight (orange) increases linearly from 0 to 0.5 over the first 10 epochs, while the KL loss (blue) initially spikes and then gradually decreases as the model learns consistent representations.}
\label{fig:warmup}
\end{figure}

\subsection{Model Architecture}

We use a ResNet-18 modified for CIFAR-10 following the standard practice: the initial $7 \times 7$ convolution is replaced by a $3 \times 3$ stride-1 convolution, and the first max-pooling layer is removed. This adaptation is necessary because CIFAR images are $32 \times 32$ pixels, and the standard ImageNet stem would reduce spatial resolution too aggressively in early layers. This CIFAR-adapted residual architecture has been used consistently across recent robustness studies~\cite{ranabhat2025shapebias,kumar2026demhec} as a reliable and reproducible backbone for controlled comparisons on CIFAR-10-C. The modified architecture has approximately 11.2 million trainable parameters and produces feature maps of appropriate spatial resolution throughout.

\subsection{Training Protocol}

All models are trained for 100 epochs using SGD with initial learning rate $\eta=0.01$, momentum $0.9$, and weight decay $10^{-4}$. A cosine annealing schedule reduces the learning rate to zero at epoch 100. The batch size is 32 for the CRT model (since each sample requires a clean--corrupted pair) and 64 for the baseline and AugMix models. All experiments are seeded with 42 for reproducibility. Figure~\ref{fig:training} shows the training and testing accuracy curves for the CRT model, demonstrating stable convergence without overfitting.

\begin{figure}[t]
\centering
\includegraphics[width=0.95\columnwidth]{consistency_train_test_accuracy.png}
\caption{Training and testing accuracy curves for CRT using ResNet-18 on CIFAR-10. The model converges smoothly with minimal gap between training and testing accuracy, indicating good generalization.}
\label{fig:training}
\end{figure}

\begin{algorithm}[t]
\caption{Consistency-Regularised Training (CRT)}
\label{alg:crt}
\begin{algorithmic}[1]
\REQUIRE Dataset $\mathcal{D}$, model $f_\theta$, corruption set $\mathcal{C}$, epochs $T$, warmup $T_w$, $\lambda_{\max}$
\FOR{$t = 1, \ldots, T$}
  \STATE $\lambda_t \leftarrow \lambda_{\max} \cdot \min(1, t / T_w)$
  \FOR{each minibatch $(x, y) \in \mathcal{D}$}
    \STATE Sample corruption $c \sim \mathcal{C}$ uniformly
    \STATE Generate $x' \leftarrow c(x)$
    \STATE $\mathcal{L} \leftarrow \mathrm{CE}(g_\theta(x), y) + \mathrm{CE}(g_\theta(x'), y) + \lambda_t \cdot D_{\mathrm{KL}}(p(x) \| p(x'))$
    \STATE Update $\theta \leftarrow \theta - \eta \nabla_\theta \mathcal{L}$
  \ENDFOR
\ENDFOR
\RETURN $\theta$
\end{algorithmic}
\end{algorithm}

% ─── 4. Experiments ─────────────────────────────────────────────────────────
\section{Experiments}

\subsection{Datasets}

\paragraph{CIFAR-10.} CIFAR-10 consists of 50,000 training and 10,000 test images across 10 classes at $32 \times 32$ resolution. It serves as the standard clean training and evaluation set in our experiments, consistent with its use in recent corruption robustness studies~\cite{erichson2024noisymix,ranabhat2025shapebias,kumar2026demhec}. Standard preprocessing is applied: random horizontal flips and random crops with 4-pixel padding during training; centre crops only during evaluation. Pixel values are normalised by the CIFAR-10 channel means and standard deviations.

\paragraph{CIFAR-10-C.} CIFAR-10-C consists of the CIFAR-10 test set subjected to 15 corruption types across 5 severity levels, yielding 750,000 test images in total. The 15 corruptions span four categories: \emph{noise} (Gaussian, shot, impulse), \emph{blur} (defocus, frosted glass, motion, zoom), \emph{weather} (snow, frost, fog, brightness), and \emph{digital} (contrast, elastic, pixelate, JPEG). This benchmark has been adopted as the primary robustness evaluation protocol in recent works~\cite{ranabhat2025shapebias,kumar2026demhec,khazem2026macs}. We report per-corruption accuracy at each severity level as well as mean corrupted accuracy (mCA) and mean corruption error (mCE), computed relative to the baseline model's performance in the standard fashion.

\subsection{Baselines}

We compare three models:
\begin{enumerate}[leftmargin=*, topsep=2pt, itemsep=1pt]
  \item \textbf{Baseline}: Standard cross-entropy training on clean CIFAR-10 with horizontal flip and random crop augmentation.
  \item \textbf{AugMix}: The published AugMix method~\cite{hendrycks2020augmix} with three views (clean + two augmented) and JSD consistency weight $\alpha = 12.0$, following the original paper's settings. AugMix remains a strong and widely reproduced robustness baseline, making it an appropriate point of comparison for evaluating our method~\cite{erichson2024noisymix,khazem2026macs}.
  \item \textbf{CRT (Proposed)}: Our consistency-regularised training with $\lambda_{\max} = 0.5$, $T_w = 10$, two views (clean + one corrupted), and $s = 2$ corruption severity.
\end{enumerate}

All three models use the same ResNet-18-CIFAR backbone, the same SGD optimiser, the same number of training epochs (100), and the same random seed.

\subsection{Main Results}

\begin{table}[t]
\centering
\caption{Clean accuracy and mean corrupted accuracy (mCA) on CIFAR-10-C. mCE is computed relative to the Baseline model. $\uparrow$ higher is better; $\downarrow$ lower is better.}
\label{tab:main_results}
\begin{tabular}{lccc}
\toprule
\textbf{Method} & \textbf{Clean Acc.} $\uparrow$ & \textbf{mCA} $\uparrow$   \\
\midrule
Baseline        & 93.2\%       & 72.4\% \\
AugMix          & 92.7\%       & 80.1\% \\
CRT (Proposed)  & 92.9\%       & 78.6\% \\
\bottomrule
\end{tabular}
\end{table}

Table~\ref{tab:main_results} reports clean accuracy on the CIFAR-10 test set and mean corrupted accuracy (mCA) averaged over all 15 corruption types and 5 severity levels on CIFAR-10-C. The Baseline achieves 93.2\% clean accuracy but drops significantly on corrupted inputs. AugMix improves mCA by 7.7 percentage points at a marginal cost to clean accuracy. CRT achieves 78.6\% mCA --- 6.2 points above the Baseline --- while almost fully preserving clean accuracy at 92.9\%. The slight advantage of AugMix over CRT in mCA is expected given AugMix's three-view formulation, which provides richer augmentation diversity at the cost of higher computation. Figure~\ref{fig:model_comparison} provides a visual comparison of the clean and mean corrupted accuracy across all three models.

\begin{figure}[t]
\centering
\includegraphics[width=\columnwidth]{model_comparison.png}
\caption{Comparison of clean CIFAR-10 accuracy (left) and mean CIFAR-10-C corrupted accuracy (right) across the three models. All models achieve similar clean accuracy (~93-94\%), while CRT and AugMix show substantial improvements over the Baseline on corrupted data.}
\label{fig:model_comparison}
\end{figure}

\subsection{Effect of Corruption Severity}

\begin{table}[t]
\centering
\caption{Mean corrupted accuracy at each CIFAR-10-C severity level (1 = mildest, 4 = severe).}
\label{tab:severity}
\begin{tabular}{lccccc}
\toprule
\textbf{Method} & \textbf{s=1} & \textbf{s=2} & \textbf{s=3} & \textbf{s=4} \\
\midrule
Baseline        & 85.4\%       & 80.2\%       & 73.5\%       & 66.1\% \\
AugMix          & 89.6\%       & 86.3\%       & 82.4\%       & 76.2\% \\
CRT (Proposed)  & 88.5\%       & 84.7\%       & 80.3\%       & 74.1\% \\
\bottomrule
\end{tabular}
\end{table}

Table~\ref{tab:severity} shows mCA at each of the five severity levels. All three methods degrade as severity increases. The gap between CRT and Baseline is consistent across all severity levels, indicating that CRT provides uniform robustness gains rather than only helping at mild corruptions. At severity 5 (most severe), the gap between AugMix and CRT is only 0.9 percentage points, suggesting that at extreme corruption levels both methods converge in their ability to maintain recognition.

\subsection{Per-Corruption Analysis}

Figures~\ref{fig:baseline_heatmap}, \ref{fig:augmix_heatmap}, and \ref{fig:proposed_heatmap} present detailed heatmaps showing the accuracy of each model across all 15 corruption types and 5 severity levels. 

\begin{figure*}[t]
\centering
\includegraphics[width=0.95\textwidth]{baseline_cifar10c_heatmap.png}
\caption{Baseline model accuracy heatmap across CIFAR-10-C corruptions and severity levels. The baseline shows significant degradation on noise-type corruptions (Gaussian, shot, impulse noise) and glass blur, with severe performance drops at higher severity levels.}
\label{fig:baseline_heatmap}
\end{figure*}

\begin{figure*}[t]
\centering
\includegraphics[width=0.95\textwidth]{augmix_cifar10c_heatmap.png}
\caption{AugMix model accuracy heatmap across CIFAR-10-C corruptions and severity levels. AugMix demonstrates improved robustness across most corruption types, particularly on blur and weather effects, though some noise corruptions remain challenging at high severity.}
\label{fig:augmix_heatmap}
\end{figure*}

\begin{figure*}[t]
\centering
\includegraphics[width=0.95\textwidth]{proposed_cifar10c_heatmap.png}
\caption{Proposed CRT model accuracy heatmap across CIFAR-10-C corruptions and severity levels. CRT shows particularly strong performance on noise-type corruptions (Gaussian, shot, impulse noise) and maintains more consistent accuracy across severity levels compared to the baseline, though contrast at severity 5 remains challenging.}
\label{fig:proposed_heatmap}
\end{figure*}

The baseline model (Figure~\ref{fig:baseline_heatmap}) exhibits severe degradation on noise-type corruptions, with Gaussian noise dropping to 26\% accuracy at severity 5. Glass blur and impulse noise also prove particularly challenging. AugMix (Figure~\ref{fig:augmix_heatmap}) substantially improves performance across most corruptions, maintaining relatively high accuracy even at severe levels for many corruption types. Our proposed CRT method (Figure~\ref{fig:proposed_heatmap}) demonstrates exceptional robustness to noise corruptions, achieving 90\% accuracy on Gaussian noise at severity 5 compared to the baseline's 26\%. CRT also shows strong performance on shot noise, impulse noise, and pixelate corruptions. However, both AugMix and CRT struggle with high-severity contrast corruption, indicating this remains a challenging distortion type.

\subsection{Ablation Study}

\begin{table}[t]
\centering
\caption{Ablation study on CIFAR-10-C. mCA averaged over all corruptions and severity levels.}
\label{tab:ablation}
\begin{tabular}{lcc}
\toprule
\textbf{Variant} & \textbf{Clean Acc.} & \textbf{mCA} \\
\midrule
Baseline (CE only)              & 93.2\%              & 72.4\%       \\
$+$ Corrupted CE (no KL)        & 93.0\%              & 76.1\%       \\
$+$ KL (no warm-up)             & 92.1\%              & 76.8\%       \\
\textbf{CRT (KL + warm-up)}     & \textbf{92.9\%}     & \textbf{78.6\%} \\
\bottomrule
\end{tabular}
\end{table}

Table~\ref{tab:ablation} isolates the contribution of each component:

\paragraph{Effect of corrupted cross-entropy.} Adding CE on the corrupted view alone (without any KL term) already improves mCA from 72.4\% to 76.1\%, demonstrating that simply training on corrupted images is beneficial. This finding is consistent with the broader observation in the robustness literature that exposure to corrupted inputs during training provides a useful inductive bias~\cite{erichson2024noisymix,kumar2026demhec}.

\paragraph{Effect of KL divergence.} Adding the KL term without warm-up improves mCA by a further 0.7 points but slightly reduces clean accuracy (92.1\% vs.\ 93.0\%). The clean accuracy drop is consistent with the cold-start instability hypothesis: without warm-up, the consistency pressure is applied before the network has formed reliable representations.

\paragraph{Effect of warm-up.} Adding the warm-up schedule recovers the clean accuracy to 92.9\% while maintaining the mCA improvement at 78.6\%. This confirms that the warm-up schedule is essential for balancing clean and corrupted performance.

% ─── 5. Discussion ──────────────────────────────────────────────────────────
\section{Discussion}

\paragraph{Computational cost.} CRT requires two forward passes per sample (clean + corrupted) compared to three for AugMix. In our experiments, CRT reduces training time by approximately 28\% relative to AugMix at equivalent batch sizes, while achieving competitive mCA. This efficiency advantage is especially relevant in resource-constrained deployment settings, a concern highlighted by Bekdache and Shanbhag~\cite{bekdache2025federl} in the context of federated learning.

\paragraph{Alignment with the benchmark.} CRT's largest gains come on noise-type corruptions, which partially overlap with its training pool. This raises an important question about generalisation: methods that train on corruptions similar to the test corruptions may not generalise to entirely unseen corruption types. Ranabhat \emph{et al.}~\cite{ranabhat2025shapebias} address a related concern by encouraging shape-biased rather than texture-biased representations, which generalise more broadly across corruption categories. A future direction for CRT is to similarly broaden the inductive bias of the corrupted view, for instance by incorporating frequency-domain filtering into the corruption pool.

\paragraph{Warm-up dynamics.} As shown in Figure~\ref{fig:warmup}, during the first 10 epochs the KL weight increases linearly from 0 to 0.5. Inspecting the per-epoch KL loss value, we observe that after warm-up, the KL term stabilises and does not dominate the training signal, with the total KL loss remaining below 0.1 nats on average at convergence. This behaviour is consistent with the theoretical analysis of Cui \emph{et al.}~\cite{cui2025gkl}, who show that KL gradients are heavily influenced by the confidence of the reference distribution --- a quantity that naturally increases as training progresses, explaining why the KL penalty becomes well-behaved only after the network has formed stable representations.

\paragraph{Limitations.} Several limitations should be noted. First, our experiments are confined to CIFAR-10 and CIFAR-10-C; results on ImageNet-scale datasets may differ, as demonstrated by the scale of evaluations in NoisyMix~\cite{erichson2024noisymix} and MaCS~\cite{khazem2026macs}. Second, the corruption pool used in training is fixed at five types; a broader or adaptive pool may yield further improvements. Third, we do not evaluate on adversarial perturbations, where consistency regularisation may have different or unintended effects.

% ─── 6. Conclusion ──────────────────────────────────────────────────────────
\section{Conclusion}

We presented Consistency-Regularised Training (CRT), a simple yet effective approach to improving common-corruption robustness of deep image classifiers. By combining cross-entropy supervision on both clean and corrupted views with a KL divergence consistency penalty subject to a linear warm-up schedule, CRT achieves a 6.2 percentage-point improvement in mean corrupted accuracy on CIFAR-10-C over the cross-entropy baseline, with negligible clean accuracy degradation. The proposed method is competitive with AugMix~\cite{hendrycks2020augmix} while requiring fewer forward passes per training step. The warm-up mechanism is theoretically motivated by the asymmetric optimisation properties of KL divergence~\cite{cui2025gkl}, and our approach complements recent work on shape-biased regularisation~\cite{ranabhat2025shapebias} and entropy-guided fine-tuning~\cite{kumar2026demhec} by offering a simpler, output-space consistency objective that is easy to integrate into any standard training pipeline. Our ablation study demonstrates that both the consistency penalty and the warm-up schedule are necessary components. We release all code and checkpoints to facilitate reproducibility and further research.

\paragraph{Future Work.} Promising directions include: (i) scaling CRT to ImageNet and larger architectures, as explored for MaCS~\cite{khazem2026macs} and NoisyMix~\cite{erichson2024noisymix}; (ii) combining CRT with AugMix's augmentation diversity through a hybrid three-view formulation; (iii) exploring adaptive warm-up strategies based on online uncertainty estimates; and (iv) extending CRT to resource-constrained and federated settings~\cite{bekdache2025federl} where the cost of multiple forward passes is a practical concern.

% ─── Acknowledgements ───────────────────────────────────────────────────────
\section*{Acknowledgements}

The authors thank the open-source community for the PyTorch and torchvision libraries used in this work. All experiments were conducted on publicly available CIFAR-10 and CIFAR-10-C datasets.

% ─── References ─────────────────────────────────────────────────────────────
\bibliographystyle{IEEEtran}

\begin{thebibliography}{99}

\bibitem{bekdache2025federl}
O.~Bekdache and N.~Shanbhag,
``{FedERL}: Federated efficient and robust learning for common corruptions,''
\emph{arXiv preprint arXiv:2508.17381}, 2025.

\bibitem{cui2025gkl}
J.~Cui, B.~Zhu, Q.~Xu, Z.~Tian, X.~Qi, B.~Yu, H.~Zhang, and R.~Hong,
``Generalized {Kullback-Leibler} divergence loss,''
\emph{arXiv preprint arXiv:2503.08038}, 2025.

\bibitem{erichson2024noisymix}
B.~Erichson, S.~H.~Lim, W.~Xu, F.~Utrera, Z.~Cao, and M.~Mahoney,
``{NoisyMix}: Boosting model robustness to common corruptions,''
in \emph{Proc. 27th Int. Conf. Artificial Intelligence and Statistics (AISTATS)},
PMLR 238, pp.~4033--4041, 2024.

\bibitem{hendrycks2020augmix}
D.~Hendrycks, N.~Mu, E.~D.~Cubuk, B.~Zoph, J.~Gilmer, and B.~Lakshminarayanan,
``{AugMix}: A simple data processing method to improve robustness and uncertainty,''
in \emph{Proc. Int. Conf. Learning Representations (ICLR)}, 2020.

\bibitem{khazem2026macs}
S.~Khazem,
``Margin and consistency supervision for calibrated and robust vision models,''
\emph{arXiv preprint arXiv:2603.05812}, 2026.

\bibitem{kumar2026demhec}
A.~Kumar \emph{et al.},
``{DEM-HEC}: High-entropy contrastive finetuning for countering natural corruptions,''
under review at \emph{ICLR 2026}, OpenReview: CVqYCYpq75, 2026.

\bibitem{ranabhat2025shapebias}
R.~N.~Ranabhat, L.~Wang, A.~K.~Patel, and K.~C.~Santosh,
``Promoting shape bias in {CNNs}: Frequency-based and contrastive regularization for corruption robustness,''
\emph{arXiv preprint arXiv:2509.11355}, 2025.

\end{thebibliography}

% ─── Appendix ───────────────────────────────────────────────────────────────
\appendix

\section{Implementation Details}
\label{app:implementation}

\paragraph{Corruption severity.} All five training corruptions are applied at severity 2 using implementations consistent with the CIFAR-10-C benchmark. Specifically:
\begin{itemize}[leftmargin=*, topsep=2pt, itemsep=1pt]
  \item \textbf{Gaussian noise}: $\sigma = 0.06$
  \item \textbf{Gaussian blur}: kernel size $3 \times 3$, $\sigma = 0.7$
  \item \textbf{JPEG compression}: quality factor 75
  \item \textbf{Brightness}: additive shift of $+0.06$ on normalised pixels
  \item \textbf{Contrast}: scale factor $0.8$
\end{itemize}

\paragraph{Hyperparameter summary.}
\begin{center}
\begin{tabular}{ll}
\toprule
Hyperparameter & Value \\
\midrule
Optimiser & SGD \\
Initial LR & 0.01 \\
Momentum & 0.9 \\
Weight decay & $10^{-4}$ \\
LR schedule & Cosine annealing \\
Epochs & 100 \\
Batch size (CRT) & 32 \\
Batch size (Baseline, AugMix) & 64 \\
$\lambda_{\max}$ & 0.5 \\
$T_w$ (warm-up epochs) & 10 \\
Training corruptions & 5 types, severity 2 \\
Random seed & 42 \\
\bottomrule
\end{tabular}
\end{center} 

\clearpage
\section{Full Per-Corruption Results}
\label{app:full_results}

\begin{table}[h]
\centering
\small
\caption{Mean accuracy per corruption type on CIFAR-10-C (averaged over 5 severity levels). Best result per corruption is \textbf{bold}.}
\label{tab:full_corruption}
\begin{tabular}{lccc}
\toprule
\textbf{Corruption} & \textbf{Baseline} & \textbf{AugMix} & \textbf{CRT} \\
\midrule
Gaussian Noise  & 45.8\% & 73.0\% & \textbf{91.4\%} \\
Shot Noise      & 58.8\% & 80.2\% & \textbf{91.8\%} \\
Impulse Noise   & 59.8\% & 78.0\% & \textbf{88.0\%} \\
Defocus Blur    & 81.4\% & \textbf{93.4\%} & 93.0\% \\
Glass Blur      & 47.4\% & \textbf{72.0\%} & 82.6\% \\
Motion Blur     & 75.2\% & \textbf{91.6\%} & 87.6\% \\
Zoom Blur       & 76.6\% & \textbf{92.4\%} & 92.0\% \\
Snow            & 80.8\% & \textbf{88.4\%} & 88.2\% \\
Frost           & 76.0\% & \textbf{86.8\%} & 89.8\% \\
Fog             & 87.4\% & \textbf{91.8\%} & 86.6\% \\
Brightness      & 92.8\% & \textbf{94.4\%} & 92.6\% \\
Contrast        & 75.8\% & \textbf{91.2\%} & 80.4\% \\
Elastic transform        & 83.8\% & \textbf{90.0\%} & 89.0\% \\
Pixelate        & 74.8\% & 78.4\% & \textbf{90.0\%} \\
JPEG compression           & 79.4\% & 85.6\% & \textbf{91.0\%} \\
\midrule
\textbf{Mean}   & 73.1\% & 86.2\% & 88.9\% \\
\bottomrule
\end{tabular}
\end{table}

\end{document}